\documentclass[11pt,a4paper]{article}
\usepackage[utf8]{inputenc}
\usepackage[T1]{fontenc}
\usepackage{lmodern,amsmath,amssymb,textcomp}
\usepackage[margin=24mm]{geometry}
\usepackage{longtable,booktabs,array,enumitem}
\usepackage[hidelinks]{hyperref}
\setlength{\parindent}{0pt}
\setlength{\parskip}{6pt}
\emergencystretch=3em
\begin{document}
{\LARGE\bfseries A three-prime exclusion for the quadratic-residue divisibility criterion\par}\medskip
{\small Provisional mathematical authors: Rachel Teo, Wenyi Wang, Hamdi Abderrahmene\par}\medskip
{\small Judging and evaluation record: the submissions discussed in this document have been judged and evaluated by Alejandro Zarzuelo Urdiales for OpenMath 2026. This dated review records the stated verification, classification and unresolved holds; it is a preliminary evaluation rather than a signed award decision.\par}

{\small Event provenance: OpenMath 2026 ran 27 September-2 October 2026 with worldwide online participation. Detailed organizer listings specify Harvard Innovation Labs for the 27 September in-person event; the OpenMath programme advertises a hybrid kickoff at MIT. Sundai identifies its Harvard/MIT community. Organizer sources: \url{https://rsihouse.ai/openmath} ; \url{https://luma.com/yzp9abvr} ; \url{https://www.sundai.club/events/boston/recursive-learning-hack-with-harvard-innovation-labs}\par}

{\small Rachel Teo  -  Sakana AI.\par}

{\small Wenyi Wang  -  Sakana AI.\par}

{\small Hamdi Abderrahmene  -  Sakana AI.\par}

{\small Affiliations follow the recorded author declarations or public institutional/profile sources. Preferred author names, author order and publication affiliations await author review. This editorial compilation does not establish anthology coauthorship.\par}

{\small Original-result working manuscript | OpenMath 2026 | 5 October 2026\par}\medskip
{\small Central source and adjudication archive: \url{https://github.com/alejandrozu/openmath-2026-judging}\par}

\subsection{Abstract}
For n=pqr with three distinct odd primes, the quadratic-residue count a(n), including zero, satisfies a(n)(a(n)\ensuremath{-}1) not dividing n\ensuremath{^2}\ensuremath{-}1. A global arithmetic reduction bounds the first two primes and the quotient; exact discriminant certificates exclude all remaining cases.

\subsection{The residue convention and theorem}
Let a(n) count all square residue classes modulo n, including zero. For distinct odd primes 2<p<q<r, the Chinese remainder theorem gives a(pqr)=(p+1)(q+1)(r+1)/8. The theorem excludes divisibility of n\ensuremath{^2}\ensuremath{-}1 by a(n)(a(n)\ensuremath{-}1) throughout this unbounded squarefree three-prime family. It therefore establishes the corresponding conjectured equivalence on this entire factorization class.

Including zero is essential: a count of only nonzero quadratic residues would change the arithmetic. The theorem says nothing about repeated prime powers or products of four or more distinct primes, and should be described as a uniform infinite subcase of A000224 rather than the full conjecture.

\subsection{Global reduction to a finite certificate}
Set A=(p+1)/2, B=(q+1)/2 and C=(r+1)/2. Under a contrary divisibility assumption, there is a positive integer quotient K satisfying (pqr)\ensuremath{^2}=KABC(ABC\ensuremath{-}1)+1. The proof derives necessary bounds from that exact equation, after explicitly justifying every conversion from natural-number subtraction to ring subtraction.

The key global step bounds p by 381 and gives a p-dependent upper bound on q. With p and q fixed, K lies in an explicit finite interval. The remaining equation is quadratic in C. Writing H=pq and U=AB, its coefficients are 4H\ensuremath{^2}\ensuremath{-}KU\ensuremath{^2}, \ensuremath{-}(4H\ensuremath{^2}\ensuremath{-}KU), and H\ensuremath{^2}\ensuremath{-}1. An integral root forces the corresponding discriminant to be an integer square.

The certificates reject that finite quotient/discriminant domain using strict gaps between consecutive squares, plus a separate congruence exclusion for the exceptional square case. Because the bounds were derived without an upper bound on r, this closes an unbounded theorem. A bounded Pell scan up to a large numerical cutoff would not by itself do so.

\subsection{Formal certificate architecture}
The source splits the argument into square-count transport, necessary bounds, quadratic reduction and generated certificate blocks. Its 171-module record includes each intended block and three final canonical endpoints. A separate Lean 4.33.1 replay is recorded in addition to 4.34.1; tactic-only compatibility changes in Bounds preserve definitions and theorem statements. Several early source comments still say unbuilt draft, so the dated manifest and selected endpoint receipts should govern verification descriptions.

For publication, append a compact definition of the finite domain and the discriminant identity, with machine-readable certificate data in the repository. The long generated leaves belong in supplementary material. Report the fresh common replay status separately from archived receipts and from reviewer sampling. An independent check of 2,024 small prime triples agrees with the formula and exclusion, but is not the universal proof.

\subsection{Prior scope and research contribution}
CRT, quotient/Pell methods and gcd obstructions are background. Earlier records already cover simpler factorization classes and a large bounded search; the claimed delta is the global reduction and exclusion for all ordered triples of distinct odd primes. The previous scan's operational n cutoff must not be confused with an unbounded theorem merely because every local survivor happens to fall below it.

This is a focused computational number-theory paper candidate if the closest pre-event project and specifications do not already entail the theorem. It should state the factorization limitation prominently and credit the baseline methods. The number of generated modules or rejected arithmetic leaves is not a count of separate discoveries.

{\small This short manuscript records the construction and selected endpoints. The companion Original results anthology contains the extended ordinary proof exposition and its explicitly stated premises; the preserved Lean source remains the complete formal proof record.\par}

\subsection{Verification and reproducibility}
Dated partial fresh replay: the 5 October 2026, 08:56:47 UTC qualification verifies 167 of 171 frozen source modules, comprising 166 granular sources and Bounds. The retained Bounds source passed under its original policy before the controller exited 120; its unfinished receipt is preserved. The subsequent source-free initial-gate failure also remains recorded. Four whole-import sources and the three original selected endpoint prints remain unrun, so the full 171-source/3-print scope is unqualified. The mathematical claim remains restricted to three ordered distinct odd primes, using the zero-inclusive quadratic-residue count; it is not the unrestricted A000224 conjecture. Qualification: Verification/A000-actual167-abnormal-controller-exit120-partial-readonly-qualification-20261005T085647463194Z.json. Historical author/internal replay is separate, and these partial checks confer no novelty, score, award or release upgrade. A later bounded dispatch assessment ended without a new Lean child or source/audit attempt. Dated resource-only summary: Verification/root-A000-whole8GiB-bounded60-public-safe-scope-20261005.json.

\begin{samepage}
\subsection{OeisA224.not\_divisible\_three\_primes}
\begin{quote}\small\ttfamily theorem not\_divisible\_three\_primes \{p q r : \ensuremath{\mathbb{N}}\}\newline     (hp : p.Prime) (hq : q.Prime) (hr : r.Prime)\newline     (hp2 : 2 < p) (hpq : p < q) (hqr : q < r) :\newline     \ensuremath{\neg} a (p * q * r) * (a (p * q * r) - 1) \ensuremath{\mid} (p * q * r) \textasciicircum{} 2 - 1\end{quote}
{\small Source: proofs/a000224/OpHack/QuadraticResidue224/Target224.lean:33.\par}

{\small Source SHA-256: 3b67f7494439886af3359a9e684433f49a6dccb7d10258d1d05dcc01de7886d0\par}

\end{samepage}
\begin{samepage}
\subsection{OeisA224.not\_modEq\_three\_primes}
\begin{quote}\small\ttfamily theorem not\_modEq\_three\_primes \{p q r : \ensuremath{\mathbb{N}}\}\newline     (hp : p.Prime) (hq : q.Prime) (hr : r.Prime)\newline     (hp2 : 2 < p) (hpq : p < q) (hqr : q < r) :\newline     \ensuremath{\neg} (p * q * r) \textasciicircum{} 2 \ensuremath{\equiv} 1 [MOD a (p * q * r) * (a (p * q * r) - 1)]\end{quote}
{\small Source: proofs/a000224/OpHack/QuadraticResidue224/Target224.lean:43.\par}

{\small Source SHA-256: 3b67f7494439886af3359a9e684433f49a6dccb7d10258d1d05dcc01de7886d0\par}

{\small\textbf{Sources and attribution:} \href{https://oeis.org/A000224}{1. oeis.org / A000224}; \href{https://github.com/umaia1234/agentic-conjectures/blob/96ea0039d9dc3a1940b9d5ec1cbe616ed58d1da2/problems/oeis-a000224/DETAILS.md}{2. umaia1234/agentic-conjectures / DETAILS.md}; \href{https://github.com/umaia1234/agentic-conjectures/blob/96ea0039d9dc3a1940b9d5ec1cbe616ed58d1da2/problems/oeis-a000224/pell_quotient_scan.py}{3. umaia1234/agentic-conjectures / pell\_quotient\_scan.py}; \href{https://github.com/alejandrozu/openmath-2026-judging/blob/d0ed487f487fa7ec814ff705ab55249983fe8707/OpenMath-Judging/review/entries/E02-a000224.txt}{4. alejandrozu/openmath-2026-judging / E02-a 000224.txt}; \href{https://github.com/alejandrozu/openmath-2026-judging}{Central source and adjudication archive}.\par}

\end{samepage}
\end{document}
